\documentclass[11pt]{article}

\usepackage[margin=1in]{geometry}
\usepackage{amsmath,amssymb,bm}
\usepackage{booktabs}
\usepackage[hidelinks]{hyperref}
\usepackage{microtype}

\newcommand{\E}{\mathbb{E}}
\newcommand{\R}{\mathbb{R}}
\newcommand{\Var}{\operatorname{Var}}
\newcommand{\Cov}{\operatorname{Cov}}
\newcommand{\tr}{\operatorname{tr}}
\newcommand{\bhat}{\widehat{\bm\beta}}
\newcommand{\bstar}{\bm\beta^\star}
\newcommand{\Shat}{\widehat\Sigma}

\title{Second-Order Fluctuations of the Accentuation Alignment Ratio\\
Current Noise-Conditional Correction and the Full Random-Design Problem}
\author{Theory supplement for \texttt{AccentuationPredRMT}}
\date{\today}

\begin{document}
\maketitle

\begin{abstract}
The own-path accentuation metrics depend on the nonlinear ratio
\[
R=\frac{\bhat^\top\bstar}{\bhat^\top\bhat}.
\]
Leading deterministic equivalents replace this random ratio by a ratio of
deterministic first and second moments.  The current code adds an approximate
second-order correction for response-noise fluctuations.  This note derives
that correction from a conditional Gaussian calculation and shows exactly
where its spectral sums arise.  It then formulates the full unconditional
variance problem, which must additionally retain random-design resolvent
fluctuations and, for cross-validated ridge, randomness of the selected
regularization.
\end{abstract}

\section{Ridge estimator and accentuation ratio}

Consider
\begin{equation}
  \bm x_i\sim\mathcal N(0,\Sigma),\qquad
  y_i=(\bstar)^\top\bm x_i+\sigma\epsilon_i,
  \qquad \epsilon_i\sim\mathcal N(0,1).
\end{equation}
For centered data, let
\begin{equation}
  \Shat=\frac{X^\top X}{n},\qquad
  Q_X=(\Shat+\lambda I)^{-1}.
\end{equation}
The normalized-loss ridge estimator is
\begin{align}
  \bhat
    &=Q_X\frac{X^\top\bm y}{n}\\
    &=Q_X\Shat\bstar
      +\frac{\sigma}{n}Q_XX^\top\bm\epsilon.
  \label{eq:ridge-decomposition}
\end{align}
Define the sample ridge filter
\begin{equation}
  P_X=Q_X\Shat=\Shat(\Shat+\lambda I)^{-1}.
\end{equation}
Conditional on \(X\), Eq.~\eqref{eq:ridge-decomposition} becomes
\begin{equation}
  \bhat=\bm\mu_X+\bm\xi_X,
  \qquad
  \bm\mu_X=P_X\bstar,
  \label{eq:conditional-decomposition}
\end{equation}
with
\begin{equation}
  \bm\xi_X\mid X\sim\mathcal N(0,\Omega_X),\qquad
  \Omega_X=\frac{\sigma^2}{n}Q_X\Shat Q_X.
  \label{eq:conditional-covariance}
\end{equation}

For one fitted model, the true-on-predicted slope along its own pixel-gradient
path is
\begin{equation}
  R=\frac{N}{D},\qquad
  N=\bhat^\top\bstar,\qquad
  D=\bhat^\top\bhat.
  \label{eq:R-definition}
\end{equation}
The normalized accentuation error and pathwise coefficient of determination
are
\begin{align}
  \frac{E_{\rm acc}}{S}&=(1-R)^2,
  \qquad S=(\bstar)^\top\Sigma\bstar,
  \label{eq:Eacc-R}\\
  R^2_{\rm acc}&=g(R)
  =1-\left(1-\frac{1}{R}\right)^2.
  \label{eq:R2acc-R}
\end{align}

\section{Exact Gaussian moments conditional on the design}

Using Eq.~\eqref{eq:conditional-decomposition},
\begin{align}
  N&=(\bstar)^\top\bm\mu_X+(\bstar)^\top\bm\xi_X,\\
  D&=\bm\mu_X^\top\bm\mu_X
  +2\bm\mu_X^\top\bm\xi_X+\bm\xi_X^\top\bm\xi_X.
\end{align}
Their conditional means are
\begin{align}
  n_X&=\E_\epsilon[N\mid X]
       =(\bstar)^\top\bm\mu_X,
  \label{eq:nX}\\
  d_X&=\E_\epsilon[D\mid X]
       =\bm\mu_X^\top\bm\mu_X+\tr\Omega_X.
  \label{eq:dX}
\end{align}
The centered fluctuations are
\begin{align}
  \delta N&=(\bstar)^\top\bm\xi_X,\\
  \delta D&=2\bm\mu_X^\top\bm\xi_X
       +\left(\bm\xi_X^\top\bm\xi_X-\tr\Omega_X\right).
\end{align}

For a centered Gaussian vector \(\bm\xi\sim\mathcal N(0,\Omega)\),
\begin{align}
  \Var(\bm a^\top\bm\xi)&=\bm a^\top\Omega\bm a,\\
  \Var(\bm\xi^\top\bm\xi)&=2\tr(\Omega^2),\\
  \Cov(\bm a^\top\bm\xi,\bm\xi^\top\bm\xi)&=0.
\end{align}
The last equality follows because centered third-order Gaussian moments vanish.
Consequently,
\begin{align}
  \Var_\epsilon(N\mid X)
    &=(\bstar)^\top\Omega_X\bstar,
  \label{eq:conditional-varN}\\
  \Cov_\epsilon(N,D\mid X)
    &=2(\bstar)^\top\Omega_X\bm\mu_X,
  \label{eq:conditional-covND}\\
  \Var_\epsilon(D\mid X)
    &=4\bm\mu_X^\top\Omega_X\bm\mu_X
      +2\tr(\Omega_X^2).
  \label{eq:conditional-varD}
\end{align}
Equations~\eqref{eq:conditional-varN}--\eqref{eq:conditional-varD} are exact
conditional Gaussian identities.

\section{Delta-method variance at fixed design}

For \(f(N,D)=N/D\), the gradient at \((n_X,d_X)\) is
\begin{equation}
  \nabla f=
  \begin{pmatrix}
    d_X^{-1}\\
    -n_Xd_X^{-2}
  \end{pmatrix}
  =
  \frac{1}{d_X}
  \begin{pmatrix}
    1\\-R_X
  \end{pmatrix},
  \qquad R_X=\frac{n_X}{d_X}.
\end{equation}
The first-order delta approximation therefore gives
\begin{align}
  \Var_\epsilon(R\mid X)
  &\simeq
  \frac{
    \Var_\epsilon(N\mid X)
    -2R_X\Cov_\epsilon(N,D\mid X)
    +R_X^2\Var_\epsilon(D\mid X)
  }{d_X^2}\\
  &=
  \frac{
    \bm v_X^\top\Omega_X\bm v_X
    +2R_X^2\tr(\Omega_X^2)
  }{d_X^2},
  \label{eq:conditional-varR}\\
  \bm v_X&=\bstar-2R_X\bm\mu_X.
\end{align}
This compact expression is the source of the current correction.  Its first
term is the variance of a linear Gaussian form; its second is the variance of
the centered quadratic form in the denominator.

\section{Deterministic substitutions made by the current code}
\label{sec:current-code}

Let
\begin{equation}
  \Sigma=U\,\mathrm{diag}(s_1,\ldots,s_d)U^\top,\qquad
  b_k=\bm u_k^\top\bstar.
\end{equation}
The effective regularization \(\kappa=\kappa(\lambda)\) solves
\begin{equation}
  \kappa-\lambda
  =\frac{\kappa}{n}\sum_k\frac{s_k}{s_k+\kappa}.
  \label{eq:kappa}
\end{equation}
Define
\begin{align}
  t_k&=\frac{s_k}{s_k+\kappa},\\
  \mathrm{df}_2&=\sum_k\frac{s_k^2}{(s_k+\kappa)^2},\\
  C_{\rm sig}&=\sum_k\frac{s_kb_k^2}{(s_k+\kappa)^2}.
\end{align}
The first-order deterministic numerator and denominator used by the code are
\begin{align}
  N_{\rm det}
    &=\sum_k t_kb_k^2,
  \label{eq:Ndet}\\
  D_{\rm det}
    &=\sum_k t_k^2b_k^2
    +\frac{\kappa^2C_{\rm sig}+\sigma^2}
           {n-\mathrm{df}_2}
       \sum_k\frac{s_k}{(s_k+\kappa)^2},
  \label{eq:Ddet}\\
  R_{\rm det}&=\frac{N_{\rm det}}{D_{\rm det}}.
  \label{eq:Rdet}
\end{align}
These equations are implemented by
\texttt{accentuation\_alignment} in
\path{rmt_core/ridge_theory_lib.py}.

For the variance correction, the current implementation makes the simplified
substitutions
\begin{align}
  \bm\mu_X&\rightsquigarrow
  T_\kappa\bstar,\qquad
  T_\kappa=\Sigma(\Sigma+\kappa I)^{-1},
  \label{eq:mu-substitution}\\
  \Omega_X&\rightsquigarrow
  \frac{\sigma^2}{n}H_\kappa,\qquad
  H_\kappa=\Sigma(\Sigma+\kappa I)^{-2},
  \label{eq:Omega-substitution}\\
  (R_X,d_X)&\rightsquigarrow(R_{\rm det},D_{\rm det}).
\end{align}
This is a hybrid leading-DE plus conditional-noise approximation.  It is not a
full central-limit theorem for the random design.

\subsection{Linear Gaussian contribution}

Under Eqs.~\eqref{eq:mu-substitution}--\eqref{eq:Omega-substitution},
\begin{equation}
  \bm v_{\rm det}
  =(I-2R_{\rm det}T_\kappa)\bstar.
\end{equation}
Therefore
\begin{align}
  \bm v_{\rm det}^\top
  \left(\frac{\sigma^2}{n}H_\kappa\right)
  \bm v_{\rm det}
  &=
  \frac{\sigma^2}{n}
  \sum_k
  \frac{s_kb_k^2}{(s_k+\kappa)^2}
  \left(1-2R_{\rm det}\frac{s_k}{s_k+\kappa}\right)^2\\
  &=
  \frac{\sigma^2}{n}
  \left(
    C_{\rm sig}
    -4R_{\rm det}M_1
    +4R_{\rm det}^2M_2
  \right),
  \label{eq:var-linear}
\end{align}
where
\begin{align}
  M_1&=\sum_k\frac{s_k^2b_k^2}{(s_k+\kappa)^3},\\
  M_2&=\sum_k\frac{s_k^3b_k^2}{(s_k+\kappa)^4}.
\end{align}
Thus the apparently specialized combination
\(C_{\rm sig}-4R_{\rm det}M_1+4R_{\rm det}^2M_2\) is simply the eigenbasis
expansion of a single quadratic form.

\subsection{Quadratic Gaussian contribution}

The second term in Eq.~\eqref{eq:conditional-varR} becomes
\begin{align}
  2R_{\rm det}^2
  \tr\left[
    \left(\frac{\sigma^2}{n}H_\kappa\right)^2
  \right]
  &=
  \frac{2R_{\rm det}^2\sigma^4}{n^2}
  T_{\rm quad},\\
  T_{\rm quad}
  &=\sum_k\frac{s_k^2}{(s_k+\kappa)^4}.
  \label{eq:var-quadratic}
\end{align}

\subsection{Current implemented formula}

Combining Eqs.~\eqref{eq:var-linear} and \eqref{eq:var-quadratic} gives
\begin{equation}
\boxed{
  \Var_{\rm current}(R)
  =
  \frac{
    \dfrac{\sigma^2}{n}
    \left(
      C_{\rm sig}
      -4R_{\rm det}M_1
      +4R_{\rm det}^2M_2
    \right)
    +
    \dfrac{2R_{\rm det}^2\sigma^4}{n^2}
    T_{\rm quad}
  }{D_{\rm det}^2}.
}
\label{eq:current-varR}
\end{equation}
Equation~\eqref{eq:current-varR} is implemented by
\texttt{accentuation\_var\_R} in
\path{rmt_core/ridge_theory_lib.py}.

The same conditional-noise substitution also gives the individual moments
needed for the mean of the ratio:
\begin{align}
  V_N&=\frac{\sigma^2}{n}C_{\rm sig},\\
  C_{ND}&=\frac{2\sigma^2}{n}M_1,\\
  V_D&=\frac{4\sigma^2}{n}M_2
      +\frac{2\sigma^4}{n^2}T_{\rm quad}.
\end{align}
The response-noise second-order shift is therefore
\begin{equation}
\boxed{
  b_{R,{\rm current}}
  =\E R-R_{\rm det}
  \simeq
  -\frac{C_{ND}}{D_{\rm det}^2}
  +\frac{R_{\rm det}V_D}{D_{\rm det}^2}.
}
\label{eq:current-meanR-bias}
\end{equation}
This is implemented by
\texttt{accentuation\_ratio\_moments\_theory}.  Like
Eq.~\eqref{eq:current-varR}, it omits random-design and CV-selection
fluctuations.

\section{How the current variance correction is used}

The exact bias--variance identity
\begin{equation}
  \E[(1-R)^2]
  =(1-\E R)^2+\Var(R)
\end{equation}
combined with Eqs.~\eqref{eq:current-varR} and
\eqref{eq:current-meanR-bias} motivates the second-order approximation
\begin{equation}
  \frac{\E E_{\rm acc}}{S}
  \simeq
  (1-R_{\rm det})^2
  +2(R_{\rm det}-1)b_{R,{\rm current}}
  +\Var_{\rm current}(R).
  \label{eq:Eacc-correction}
\end{equation}
The variance term is enabled by \path{include_var_R=True}; the mean-ratio
term is enabled by \path{include_ratio_mean_correction=True}.  Omitting both yields the
ratio-of-deterministic-expectations curve.

For pathwise \(R^2\), let
\begin{equation}
  g(r)=1-\left(1-\frac1r\right)^2
      =\frac{2}{r}-\frac{1}{r^2}.
\end{equation}
A second-order scalar delta expansion gives
\begin{align}
  \E[g(R)]
  &\simeq
  g(R_{\rm det})
  +\frac12g''(R_{\rm det})\Var_{\rm current}(R),\\
  \frac12g''(r)&=\frac{2r-3}{r^4}.
\end{align}
Hence
\begin{equation}
\boxed{
  R_{\rm acc}^{2,\rm corrected}
  =
  1-\left(1-\frac{1}{R_{\rm det}}\right)^2
  +\frac{2(1-R_{\rm det})}{R_{\rm det}^3}
   b_{R,{\rm current}}
  +
  \frac{2R_{\rm det}-3}{R_{\rm det}^4}
  \Var_{\rm current}(R).
}
\label{eq:R2-correction}
\end{equation}
The variance term is enabled by \path{include_delta_correction=True}, and the
mean-ratio term by \path{include_ratio_mean_correction=True}.
Equation~\eqref{eq:R2-correction} is poorly conditioned when
\(R_{\rm det}\) is close to zero; in that regime the local Taylor expansion
need not describe the heavy-tailed reciprocal statistic.

\section{What conditioning on the design omits}

Conditioning on \(X\) is an exact first step.  The incomplete step is replacing
all \(X\)-dependent conditional means by deterministic limits while retaining
only response-noise variance.  The law of total variance is
\begin{equation}
\boxed{
  \Var_{X,\epsilon}(R)
  =
  \E_X\!\left[\Var_\epsilon(R\mid X)\right]
  +
  \Var_X\!\left(\E_\epsilon[R\mid X]\right).
}
\label{eq:total-variance}
\end{equation}
The current correction targets primarily the first term.  The second term is
a random-design fluctuation of the ridge resolvent.

This distinction is invisible to a leading deterministic equivalent: both
terms can be \(O(n^{-1})\) and therefore vanish as \(n\to\infty\), while their
sum is precisely the leading quantity needed for a second-order correction.
A sharp diagnostic is \(\sigma=0\).  Equation~\eqref{eq:current-varR} then
returns zero, although finite-sample variation of \(P_X\bstar\) across
training sets generally gives nonzero \(\Var_X(R)\).

The covariance substitution in Eq.~\eqref{eq:Omega-substitution} is also
simplified.  Deterministic equivalents for derivatives of the resolvent
introduce the susceptibility
\begin{equation}
  \kappa'(\lambda)
  =
  \frac{1}{
  1-\mathrm{df}_2/n
  }.
  \label{eq:kappa-prime}
\end{equation}
A minimal refinement should carry such factors in noise bilinear forms.
However, simply multiplying every term by \(\kappa'\) is not a complete
solution: the quadratic trace and design covariance require joint
two-resolvent equivalents.

\section{Formulation of the full unconditional variance}
\label{sec:full-varR}

Let
\begin{equation}
  \bm Z=
  \begin{pmatrix}N\\D\end{pmatrix},\qquad
  \overline{\bm Z}=
  \begin{pmatrix}\overline N\\\overline D\end{pmatrix}
  =
  \E_{X,\epsilon}\bm Z.
\end{equation}
To leading second order,
\begin{equation}
  \Var(R)
  \simeq
  \bm J^\top\mathcal C_{ND}\bm J,
  \qquad
  \bm J=
  \begin{pmatrix}
    \overline D^{-1}\\
    -\overline N\overline D^{-2}
  \end{pmatrix},
  \label{eq:full-delta}
\end{equation}
where
\begin{equation}
  \mathcal C_{ND}
  =
  \Cov_{X,\epsilon}
  \begin{pmatrix}N\\D\end{pmatrix}.
\end{equation}
The exact covariance decomposition is
\begin{equation}
\boxed{
  \mathcal C_{ND}
  =
  \E_X[\mathcal C_{\epsilon}(X)]
  +
  \mathcal C_X,
}
\label{eq:covariance-decomposition}
\end{equation}
with
\begin{equation}
  \mathcal C_{\epsilon}(X)
  =
  \begin{pmatrix}
    (\bstar)^\top\Omega_X\bstar
    &
    2(\bstar)^\top\Omega_X\bm\mu_X\\
    2(\bstar)^\top\Omega_X\bm\mu_X
    &
    4\bm\mu_X^\top\Omega_X\bm\mu_X
      +2\tr(\Omega_X^2)
  \end{pmatrix},
  \label{eq:conditional-cov-matrix}
\end{equation}
and
\begin{equation}
  \mathcal C_X
  =
  \Cov_X
  \begin{pmatrix}
    n_X\\d_X
  \end{pmatrix}.
  \label{eq:design-cov-matrix}
\end{equation}
The current formula approximates selected entries of
\(\E_X[\mathcal C_{\epsilon}(X)]\) and sets
\(\mathcal C_X\) to zero.

\subsection{Random-matrix objects required for the design covariance}

Using \(P_X=\Shat(\Shat+\lambda I)^{-1}\),
\begin{align}
  n_X&=(\bstar)^\top P_X\bstar,\\
  d_X&=(\bstar)^\top P_X^2\bstar
       +\frac{\sigma^2}{n}
        \tr(Q_X\Shat Q_X).
  \label{eq:nX-dX-resolvent}
\end{align}
Since
\begin{equation}
  P_X=I-\lambda Q_X,\qquad
  P_X^2=I-2\lambda Q_X+\lambda^2Q_X^2,
\end{equation}
the design covariance in Eq.~\eqref{eq:design-cov-matrix} requires a joint
central-limit theorem for bilinear and trace resolvent statistics, including
\begin{align}
  &(\bstar)^\top Q_X\bstar,\qquad
  (\bstar)^\top Q_X^2\bstar,\\
  &\tr(Q_X\Shat Q_X),\\
  &\Cov\left(
    (\bstar)^\top Q_X\bstar,
    (\bstar)^\top Q_X^2\bstar
  \right),\\
  &\Cov\left(
    (\bstar)^\top Q_X\bstar,
    \tr(Q_X\Shat Q_X)
  \right),
\end{align}
and the remaining entries generated by Eq.~\eqref{eq:nX-dX-resolvent}.
Equivalently, these can be obtained from a two-spectral-parameter covariance
kernel for
\[
(\bstar)^\top(\Shat-zI)^{-1}\bstar
\]
and its derivatives before taking \(z=-\lambda\).

A complete calculation must specify:
\begin{enumerate}
  \item deterministic means of \(N\) and \(D\), including \(O(n^{-1})\)
        bias if the mean of \(R\) is required to the same order;
  \item the full \(2\times2\) covariance matrix
        \(\mathcal C_{ND}\) to order \(n^{-1}\);
  \item fourth-cumulant terms for non-Gaussian designs;
  \item spike or teacher-alignment terms when \(\bstar\) is not generic with
        respect to the covariance eigenbasis.
\end{enumerate}
For Gaussian designs the fourth-cumulant correction vanishes, but bilinear
resolvent fluctuations remain.

\subsection{Second-order correction to the mean ratio}

If \(\E R\), rather than only \(\Var R\), is needed to order \(n^{-1}\), the
second-order bivariate delta method gives
\begin{equation}
  \E R
  \simeq
  \frac{\overline N}{\overline D}
  -\frac{\Cov(N,D)}{\overline D^2}
  +\frac{\overline N\,\Var(D)}{\overline D^3}.
  \label{eq:mean-ratio-correction}
\end{equation}
Thus a full correction to \(\E[(1-R)^2]\) should use both the corrected mean
in Eq.~\eqref{eq:mean-ratio-correction} and the full variance in
Eq.~\eqref{eq:full-delta}.

\section{Additional randomness from cross-validation}

The fixed-\(\lambda\) theory above treats the regularization as deterministic.
For RidgeCV, the fitted model is
\begin{equation}
  \bhat=\bhat_{\widehat\alpha(X,\epsilon)},
\end{equation}
so the selected grid point is another random variable.  The unconditional
variance contains:
\begin{enumerate}
  \item fluctuations of \(N\) and \(D\) at fixed \(\alpha\);
  \item variation of \(\widehat\alpha\) across data and response noise;
  \item covariance between \(\widehat\alpha\) and the downstream ratio.
\end{enumerate}
If the CV risk has a unique smooth interior minimum, a local approximation is
\begin{equation}
  \widehat\alpha-\alpha_\star
  \simeq
  -\frac{\widehat L_{\rm CV}'(\alpha_\star)}
         {L_{\rm CV}''(\alpha_\star)}.
  \label{eq:cv-alpha-delta}
\end{equation}
One could augment the joint delta calculation with
\((N,D,\widehat\alpha)\).  For a discrete grid or a flat risk curve, the local
approximation in Eq.~\eqref{eq:cv-alpha-delta} is unreliable; the appropriate
description is a mixture over selected grid points,
\begin{equation}
  \Var(R)
  =
  \sum_a p_a\Var(R\mid\widehat\alpha=a)
  +
  \Var_a\!\left(\E[R\mid\widehat\alpha=a]\right).
\end{equation}

\section{Linear-feature generalization}

For a linear feature map, let the feature coefficient be
\(\widehat{\bm\theta}=\bm\mu+\bm\xi\), with conditional covariance \(\Omega\).
Pixel back-propagation introduces
\begin{equation}
  G=AA^\top,\qquad \bm q=A\bstar,
\end{equation}
and the own-path ratio is
\begin{equation}
  R=
  \frac{\widehat{\bm\theta}^\top\bm q}
       {\widehat{\bm\theta}^\top G\widehat{\bm\theta}}.
\end{equation}
The fixed-design delta calculation becomes
\begin{equation}
\boxed{
  \Var(R\mid X)
  \simeq
  \frac{
    \bm v^\top\Omega\bm v
    +2R^2\tr[(G\Omega)^2]
  }{D^2},
  \qquad
  \bm v=\bm q-2RG\bm\mu.
}
\label{eq:feature-varR}
\end{equation}
The pixel-space formula is recovered with \(G=I\) and \(\bm q=\bstar\).
For noncommuting feature covariance and back-propagation geometry, the full
unconditional theory requires weighted two-resolvent covariances.

\section{Practical validation and implementation roadmap}

A nested simulation can measure every component of
Eq.~\eqref{eq:total-variance} without theoretical ambiguity:
\begin{enumerate}
  \item draw \(B_X\) independent training matrices;
  \item for each fixed \(X_b\), draw \(B_\epsilon\) independent response-noise
        vectors;
  \item estimate the within-\(X\) variance and average it over datasets;
  \item estimate the between-\(X\) variance of the conditional means;
  \item repeat first at fixed \(\alpha\), then with RidgeCV.
\end{enumerate}
This separates response noise, random design, and selection variability.

A staged theoretical implementation should proceed as follows:
\begin{enumerate}
  \item replace the simplified noise covariance in
        Eq.~\eqref{eq:Omega-substitution} with its derivative-resolvent DE;
  \item derive or import the joint Gaussian-design CLT for the bilinear forms
        in Section~\ref{sec:full-varR};
  \item assemble \(\mathcal C_{ND}\) using
        Eq.~\eqref{eq:covariance-decomposition};
  \item apply Eqs.~\eqref{eq:full-delta} and
        \eqref{eq:mean-ratio-correction};
  \item validate fixed-\(\alpha\) predictions with nested Monte Carlo;
  \item add CV-selection variability only after the fixed-\(\alpha\) variance
        is validated.
\end{enumerate}

\section{Summary}

The current code formula is not an arbitrary spectral expression.  It follows
from four steps: conditional Gaussian ridge noise, a delta method for \(N/D\),
Wick identities, and a leading spectral DE substitution.
It estimates response-noise fluctuations around a deterministic ridge fit.
The full \(\Var(R)\) additionally needs the joint random-design covariance of
the numerator and denominator resolvent statistics, second-order correction
to the mean ratio, and--for RidgeCV--the distribution and covariance of the
selected regularization.

\end{document}
